% from ega4-11-fr.tex
\oldpage[IV-3]{167}

\begin{proposition}[11.9.15]
\label{IV.11.9.15-fr}
Les notations étant celles de (11.9.1), supposons que $X$ soit un
$S$-préschéma, les $\mathcal G_\lambda$ étant $S$-plats. Soit $S_0$ un
sous-préschéma fermé de $S$ défini par un Idéal quasi-cohérent nilpotent
$\mathcal J$ de $\mathcal O_S$, tel que les
$(\mathcal O_S/\mathcal J)$-Modules $\mathcal J^k/\mathcal J^{k+1}$ soient
tous localement libres. Soit $(u_{\lambda0})$ la famille d'homomorphismes
obtenue à partir du changement de base
$X\leftarrow X_0=X\times_S S_0$. Alors, si la famille $(u_{\lambda0})$ est
séparante (resp. universellement séparante relativement à $S_0$), la famille
$(u_\lambda)$ est séparante (resp. universellement séparante relativement à
$S$).
\end{proposition}

\begin{proof}
Notons que si $S'\to S$ est un changement de base quelconque et
$S'_0=S'\times_S S_0$, $S'_0$ est un sous-préschéma fermé de $S'$ défini par
un idéal quasi-cohérent nilpotent $\mathcal J'$ de $\mathcal O_{S'}$, tel que
$\mathcal J'^k/\mathcal J'^{k+1}$ soit un
$(\mathcal O_{S'}/\mathcal J')$-Module localement libre pour tout $k$
(2.1.8, (i))~; comme en outre les
$\mathcal G'_\lambda=\mathcal G_\lambda\otimes_S S'$ sont $S'$-plats, on voit
que l'assertion relative aux familles universellement séparantes est
conséquence de l'assertion relative aux familles séparantes. Pour prouver
cette dernière, on peut (11.9.3) se ramener au cas où
$S=\operatorname{Spec}(A)$, $X=\operatorname{Spec}(B)$ sont affines,
$Z_\lambda=X$ pour tout $\lambda$, $\mathcal F=\widetilde M$,
$\mathcal G_\lambda=\widetilde N_\lambda$, où $M$ et les $N_\lambda$ sont des
$B$-modules, les $N_\lambda$ étant des $A$-modules plats. En outre, la question
étant locale sur $X$ et $S$, il suffit de voir que si $t\in M$ est tel que
$u_\lambda(t)=0$ pour tout $\lambda$, alors $t=0$. On a
$\mathcal J=\widetilde{\mathfrak J}$, où $\mathfrak J$ est un idéal nilpotent
de $A$, tel que les $\mathfrak J^k/\mathfrak J^{k+1}$ soient des
$(A/\mathfrak J)$-modules libres, et par hypothèse les
$u_{\lambda0}:M/\mathfrak JM\to N_\lambda/\mathfrak JN_\lambda$ forment une
famille séparante. Supposons que $\mathfrak J^{n+1}=0$ ($n$ entier $\geq0$) et
raisonnons par récurrence sur $n$, l'assertion étant triviale pour $n=0$. Si
$\bar t$ est la classe de $t$ dans $M/\mathfrak J^nM$, la classe de
$u_\lambda(t)$ dans $N_\lambda/\mathfrak J^nN_\lambda$ est nulle pour tout
$\lambda$, donc, par l'hypothèse de récurrence, $\bar t=0$, autrement dit on a
$t\in\mathfrak J^nM$. Or
$\mathfrak J^n=\mathfrak J^n/\mathfrak J^{n+1}$ est un
$(A/\mathfrak J)$-module libre~; si $(e_\alpha)$ est une base de ce module, on
peut donc écrire $t=\sum_\alpha e_\alpha t^0_\alpha$, avec
$t^0_\alpha\in M/\mathfrak JM$, nul sauf pour un nombre fini d'indices.
D'autre part, puisque $N_\lambda$ est un $A$-module plat,
$\mathfrak J^nN_\lambda$ s'identifie à
$\mathfrak J^n\otimes_{(A/\mathfrak J)}(N_\lambda/\mathfrak JN_\lambda)$ et
l'on peut par suite écrire
\[
u_\lambda(t)=\sum_\alpha e_\alpha u_{\lambda0}(t^0_\alpha)
=\sum_\alpha e_\alpha\otimes u_{\lambda0}(t^0_\alpha).
\]
Comme par hypothèse $u_\lambda(t)=0$, on en déduit que
$u_{\lambda0}(t^0_\alpha)=0$ pour tout $\alpha$ et tout $\lambda$~; d'où
$t^0_\alpha=0$ pour tout $\alpha$ puisque la famille $(u_{\lambda0})$ est
séparante, et par suite $t=0$. C.Q.F.D.
\end{proof}

\begin{theorem}[11.9.16]
\label{IV.11.9.16-fr}
Les notations étant celles de (11.9.1), supposons que $X$ soit un
$S$-préschéma localement noethérien, $\mathcal F$ un $\mathcal O_X$-Module
cohérent et que les $\mathcal G_\lambda$ soient $S$-plats. Pour tout $s\in S$,
soit $((u_\lambda)_s)_{\lambda\in L}$ la famille obtenue à partir de
$(u_\lambda)$ par le changement de base
$X_s=X\times_S\operatorname{Spec}(k(s))\to X$. Alors, pour que la famille
$(u_\lambda)$ soit universellement séparante relativement à $S$, il faut et il
suffit que pour tout $s\in S$, la famille $((u_\lambda)_s)$ soit séparante.
\end{theorem}

\begin{proof}
La nécessité de la condition découle trivialement des définitions.
Inversement, supposons la condition de l'énoncé vérifiée, et prouvons d'abord
que la famille $(u_\lambda)$ est séparante. On peut (11.9.3) se réduire au cas
où $S=\operatorname{Spec}(A)$, $X=\operatorname{Spec}(B)$ sont affines,
$Z_\lambda=X$ pour tout $\lambda$, $\mathcal F=\widetilde M$,
$\mathcal G_\lambda=\widetilde N_\lambda$, où $B$ est noethérien, $M$ est un
$B$-module de type fini et les $N_\lambda$ des $A$-modules plats et se borner à
prouver que, si $t\in M$ est tel que $u_\lambda(t)=0$ pour tout $\lambda$,
alors $t=0$. Pour montrer que $t=0$, il suffit de prouver que pour tout idéal
maximal $\mathfrak p$ de $B$, l'image $t_\mathfrak p$ de $t$ dans
$M_\mathfrak p$ est nulle (Bourbaki, \emph{Alg. comm.}, chap. II, \S~3,
n\textsuperscript{o}~3, cor. 1 du th. 1). On peut donc se borner à montrer que
l'intersection
